\chapter{Математическая и численная модель теплопередачи}

\section{Физические основы теплопроводности в сплошных средах}

Теплопроводность является одним из основных механизмов переноса тепловой энергии в сплошных средах и обусловлена взаимодействием частиц вещества на микроскопическом уровне~\cite{mayer1933}. В твёрдых телах перенос тепла осуществляется за счёт колебаний кристаллической решётки и движения свободных электронов, что приводит к выравниванию температуры внутри материала.

Математическое описание процесса теплопроводности~\cite{сергиенко2004исследование} основывается на законе сохранения энергии, согласно которому изменение внутренней энергии элементарного объёма среды определяется балансом подводимого и отводимого тепла. В общем случае изменение температуры во времени связано с тепловыми потоками и внутренними источниками тепла.

Основным законом, описывающим процесс теплопроводности~\cite{luo2010effect}, является закон Фурье, согласно которому плотность теплового потока пропорциональна градиенту температуры и направлена в сторону её убывания.

На основе закона сохранения энергии и закона Фурье выводится уравнение теплопроводности, также известное как уравнение Фурье--Кирхгофа~\cite{chichko2013mathematical}. В трёхмерном случае оно описывает изменение температуры во времени в зависимости от пространственного распределения температуры и физических свойств материала.


\subsection{Вывод уравнения теплопроводности Фурье--Кирхгофа}

Вывод уравнения теплопроводности основывается на законе сохранения энергии, применённом к элементарному объёму сплошной среды. Рассмотрим небольшой объём материала, в котором температура является функцией координат и времени $T = T(x,y,z,t)$.

Изменение внутренней энергии данного объёма за время $dt$ определяется выражением:
\begin{equation}
\rho c \frac{\partial T}{\partial t} dV,
\end{equation}
где $\rho$~--- плотность вещества, $c$~--- удельная теплоёмкость.

Тепловой поток в среде описывается законом Фурье:
\begin{equation}
\mathbf{q} = -k \nabla T,
\end{equation}
где $\mathbf{q}$~--- вектор плотности теплового потока, $k$~--- коэффициент теплопроводности, $\nabla T$~--- градиент температуры.

Количество тепла, выходящего из объёма, определяется дивергенцией теплового потока:
\begin{equation}
\nabla \cdot \mathbf{q}.
\end{equation}

С учётом возможного наличия внутренних источников тепла с объёмной плотностью $Q$ уравнение энергетического баланса принимает вид:
\begin{equation}
\rho c \frac{\partial T}{\partial t} = - \nabla \cdot \mathbf{q} + Q.
\end{equation}

Подставляя закон Фурье в данное выражение, получаем:
\begin{equation}
\rho c \frac{\partial T}{\partial t} = \nabla \cdot (k \nabla T) + Q.
\end{equation}

Полученное уравнение является уравнением теплопроводности Фурье--Кирхгофа в общем виде.

В случае однородного изотропного материала с постоянным коэффициентом теплопроводности $k$ уравнение упрощается:
\begin{equation}
\rho c \frac{\partial T}{\partial t} = k \Delta T + Q,
\end{equation}
где $\Delta$~--- оператор Лапласа.

Введём коэффициент температуропроводности:
\begin{equation}
a = \frac{k}{\rho c},
\end{equation}
тогда уравнение принимает вид:
\begin{equation}
\frac{\partial T}{\partial t} = a \Delta T + \frac{Q}{\rho c}.
\end{equation}

Данное уравнение используется в дальнейшем для численного моделирования нестационарных процессов теплопередачи.

\subsection{Коэффициент температуропроводности и физические свойства материалов}

Коэффициент температуропроводности~\cite{vlasov2004research} $a$ является ключевым физическим 
параметром, определяющим скорость выравнивания температуры в материале 
при нестационарных процессах теплопроводности. Он связывает теплопроводность 
материала с его теплоёмкостью и плотностью:
\begin{equation}
    a = \frac{k}{\rho c},
\end{equation}
где $k$~--- коэффициент теплопроводности, $\rho$~--- 
плотность вещества, $c$~--- удельная теплоёмкость.

Физический смысл коэффициента температуропроводности состоит в следующем: 
чем выше его значение, тем быстрее тепло распространяется по материалу 
и тем быстрее устанавливается тепловое равновесие. Материалы с высоким 
значением $a$ быстро реагируют на изменение температурных условий, тогда 
как материалы с низким значением $a$ характеризуются медленным 
перераспределением тепла.

Среди металлов наибольшим коэффициентом температуропроводности обладают 
медь~($a = 1{,}17 \cdot 10^{-4}$~м$^2$/с) и алюминий~($a = 9{,}76 \cdot 
10^{-5}$~м$^2$/с), что обусловливает их широкое применение в системах 
теплоотвода. Конструкционная сталь характеризуется значением $a = 1{,}28 
\cdot 10^{-5}$~м$^2$/с и является наиболее распространённым материалом в 
инженерных расчётах. Нержавеющая сталь и титан обладают несколько меньшей 
температуропроводностью --- $3{,}91 \cdot 10^{-6}$~м$^2$/с и $9{,}40 \cdot 
10^{-6}$~м$^2$/с соответственно.

Строительные и керамические материалы характеризуются существенно более 
низкими значениями коэффициента температуропроводности. Для бетона 
$a = 6{,}42 \cdot 10^{-7}$~м$^2$/с, для гранита --- $1{,}41 \cdot 
10^{-6}$~м$^2$/с, для кирпича --- $5{,}29 \cdot 10^{-7}$~м$^2$/с. 
Это определяет их теплоизолирующие свойства и медленное перераспределение 
тепла при изменении температурных условий.

Перечисленные материалы включены в библиотеку программного комплекса 
HeatSim3D и доступны для выбора при проведении расчётов без необходимости 
ручного ввода теплофизических параметров.

\subsection{Постановка задачи численного моделирования процесса теплопроводности}

Численное моделирование нестационарной теплопроводности выполняется 
в трёхмерной прямоугольной области $\Omega = [0, L_x] \times [0, L_y] 
\times [0, L_z]$. Температура в каждой точке области является функцией 
координат и времени $T = T(x, y, z, t)$ и удовлетворяет уравнению 
теплопроводности Фурье--Кирхгофа:
\begin{equation}
    \frac{\partial T}{\partial t} = a \left( 
    \frac{\partial^2 T}{\partial x^2} + 
    \frac{\partial^2 T}{\partial y^2} + 
    \frac{\partial^2 T}{\partial z^2} \right),
\end{equation}
где $a$~--- коэффициент температуропроводности материала.

\subsection{Начальные и граничные условия трёхмерной задачи теплопроводности}

В начальный момент времени температура во всей расчётной 
области предполагается равномерной:
\begin{equation}
    T(x, y, z, 0) = T_0, \quad (x, y, z) \in \Omega,
\end{equation}
где $T_0$~--- начальная температура тела.

\subsection{Граничные условия Дирихле для шести граней области}

На всех гранях расчётной области задаются граничные условия первого 
рода, при которых температура фиксируется и поддерживается постоянной 
на протяжении всего времени моделирования. Для каждой из шести граней 
температура задаётся независимо:
\begin{equation}
    T\big|_{x=0} = T_{x^-}, \quad T\big|_{x=L_x} = T_{x^+},
\end{equation}
\begin{equation}
    T\big|_{y=0} = T_{y^-}, \quad T\big|_{y=L_y} = T_{y^+},
\end{equation}
\begin{equation}
    T\big|_{z=0} = T_{z^-}, \quad T\big|_{z=L_z} = T_{z^+}.
\end{equation}

Задание независимых температур на каждой грани позволяет моделировать 
как равномерный нагрев, так и направленный тепловой поток через тело.

\subsection{Расчётная область и основные параметры сетки}

Расчётная область разбивается на равномерную декартову сетку 
с $N_x \times N_y \times N_z$ узлами. Шаги сетки определяются как:
\begin{equation}
    \Delta x = \frac{L_x}{N_x - 1},
\end{equation}
\begin{equation}
    \Delta y = \frac{L_y}{N_y - 1},
\end{equation}
\begin{equation}
    \Delta z = \frac{L_z}{N_z - 1}.
\end{equation}

При увеличении числа узлов точность численного решения возрастает, 
однако при этом увеличиваются вычислительные затраты. В программном 
комплексе HeatSim3D параметры сетки $N_x$, $N_y$, $N_z$ задаются 
пользователем в зависимости от требуемой точности расчёта.

\section{Явная конечно-разностная схема FTCS}

Для численного решения уравнения теплопроводности применяется явная 
конечно-разностная~\cite{kolgan1975finite} схема FTCS (Forward Time, Centered Space) --- схема 
с прямым ходом по времени и центральными разностями по пространству.

Значение температуры в узле сетки с индексами $(i, j, k)$ в момент 
времени $t^n$ обозначается $T_{i,j,k}^n$. Временной шаг обозначается 
$\Delta t$, пространственные шаги --- $\Delta x$, $\Delta y$, $\Delta z$.

\subsection{Дискретизация уравнения теплопроводности по пространству и времени}

Производная по времени аппроксимируется разностью вперёд первого 
порядка точности:
\begin{equation}
    \frac{\partial T}{\partial t} \approx 
    \frac{T_{i,j,k}^{n+1} - T_{i,j,k}^{n}}{\Delta t}.
\end{equation}

Вторые производные по пространственным координатам аппроксимируются 
центральными разностями второго порядка точности:
\begin{equation}
    \frac{\partial^2 T}{\partial x^2} \approx 
    \frac{T_{i+1,j,k}^{n} - 2T_{i,j,k}^{n} + T_{i-1,j,k}^{n}}{\Delta x^2},
\end{equation}
\begin{equation}
    \frac{\partial^2 T}{\partial y^2} \approx 
    \frac{T_{i,j+1,k}^{n} - 2T_{i,j,k}^{n} + T_{i,j-1,k}^{n}}{\Delta y^2},
\end{equation}
\begin{equation}
    \frac{\partial^2 T}{\partial z^2} \approx 
    \frac{T_{i,j,k+1}^{n} - 2T_{i,j,k}^{n} + T_{i,j,k-1}^{n}}{\Delta z^2}.
\end{equation}

Подставляя данные аппроксимации в уравнение теплопроводности, получаем 
явную разностную схему для вычисления температуры на новом временном 
слое:
\begin{equation}
    T_{i,j,k}^{n+1} = T_{i,j,k}^{n} + 
    r_x \delta_x^2 T + 
    r_y \delta_y^2 T + 
    r_z \delta_z^2 T,
\end{equation}
где введены обозначения для разностных операторов:
\begin{equation}
    \delta_x^2 T = T_{i+1,j,k}^{n} - 2T_{i,j,k}^{n} + T_{i-1,j,k}^{n},
\end{equation}
\begin{equation}
    \delta_y^2 T = T_{i,j+1,k}^{n} - 2T_{i,j,k}^{n} + T_{i,j-1,k}^{n},
\end{equation}
\begin{equation}
    \delta_z^2 T = T_{i,j,k+1}^{n} - 2T_{i,j,k}^{n} + T_{i,j,k-1}^{n},
\end{equation}
и коэффициентов схемы:
\begin{equation}
    r_x = \frac{a \Delta t}{\Delta x^2}, \quad
    r_y = \frac{a \Delta t}{\Delta y^2}, \quad
    r_z = \frac{a \Delta t}{\Delta z^2}.
\end{equation}

\subsection{Условие устойчивости --- критерий Куранта--Фридрихса--Леви}

Явная схема FTCS~\cite{dubey2015stability} является условно устойчивой. Условие устойчивости 
в трёхмерном случае имеет вид:
\begin{equation}
    2a \Delta t \left( 
    \frac{1}{\Delta x^2} + 
    \frac{1}{\Delta y^2} + 
    \frac{1}{\Delta z^2} 
    \right) \leq 1.
\end{equation}

При нарушении данного условия численное решение становится неустойчивым 
и не имеет физического смысла. В программном комплексе HeatSim3D 
временной шаг $\Delta t$ выбирается автоматически из условия 
устойчивости с коэффициентом запаса $0{,}4$:
\begin{equation}
    \Delta t = \frac{0{,}4}{2a \left( 
    \dfrac{1}{\Delta x^2} + 
    \dfrac{1}{\Delta y^2} + 
    \dfrac{1}{\Delta z^2} \right)}.
\end{equation}

Это гарантирует выполнение условия устойчивости при любых параметрах 
расчётной сетки и физических свойствах материала.

\subsection{Оценка погрешности аппроксимации и сходимости разностной схемы}

Аппроксимационная погрешность схемы FTCS определяется порядками 
точности применённых разностных аппроксимаций. Погрешность по времени 
имеет первый порядок, по пространственным координатам --- 
второй порядок. Таким образом, 
суммарная погрешность схемы составляет:
\begin{equation}
    O(\Delta t) + O(\Delta x^2 + \Delta y^2 + \Delta z^2).
\end{equation}

Уменьшение шагов сетки и временного шага приводит к снижению 
погрешности и сходимости численного решения к точному. При этом 
временной шаг  должен удовлетворять условию устойчивости, 
что накладывает ограничение на минимально допустимый пространственный 
шаг при фиксированном $\Delta t$.

